\documentclass[final]{beamer}

\usepackage[size=a0,orientation=landscape,scale=1.08]{beamerposter}

\usepackage[utf8]{inputenc}
\usepackage[T1]{fontenc}
\usepackage[english]{babel}
\usepackage{amsmath,amssymb}
\usepackage{booktabs}
\usepackage{array}
\usepackage{graphicx}
\usepackage{xcolor}
\usepackage{ragged2e}
\usepackage{tikz}
\usetikzlibrary{arrows.meta,calc,decorations.pathreplacing,positioning}

\usetheme{default}
\setbeamertemplate{navigation symbols}{}
\setbeamercolor{normal text}{fg=black,bg=white}
\setbeamercolor{headline}{fg=black,bg=white}
\setbeamercolor{footline}{fg=black,bg=white}

\definecolor{titleblue}{RGB}{24,74,126}
\definecolor{rulegray}{RGB}{176,184,194}
\definecolor{textgray}{RGB}{45,50,58}
\definecolor{mutedgray}{RGB}{104,112,122}
\definecolor{panelgray}{RGB}{246,248,250}
\definecolor{peftblue}{RGB}{20,96,170}
\definecolor{doragreen}{RGB}{34,132,84}
\definecolor{proorange}{RGB}{202,98,28}
\definecolor{zeroshotgray}{RGB}{120,128,138}
\definecolor{linearblue}{RGB}{38,115,190}
\definecolor{lorapurple}{RGB}{126,87,194}

\setlength{\parindent}{0pt}
\setlength{\parskip}{0.35em}
\renewcommand{\arraystretch}{1.18}

\newcommand{\code}[1]{\texttt{#1}}
\newcommand{\method}[1]{\textbf{#1}}
\newcommand{\sectiontitle}[1]{%
  \vspace{0.45em}%
  {\Large\bfseries\color{titleblue}#1}\par
  \vspace{0.10em}%
  {\color{rulegray}\rule{\linewidth}{0.9pt}}\par
  \vspace{0.22em}%
}
\newcommand{\figureplaceholder}[1]{%
  \begin{center}
  \begin{tikzpicture}
    \draw[rulegray, line width=1pt] (0,0) rectangle (0.92\linewidth,4.5);
    \node[align=center, text=mutedgray] at (0.46\linewidth,2.25) {#1};
  \end{tikzpicture}
  \end{center}
}

\begin{document}
\begin{frame}[t]

\begin{center}
  {\VeryHuge\bfseries\color{titleblue}
  Investigating Gradient Mismatch in Low-Rank Adaptation of CLIP-ViT}\\[0.35em]
  {\Large Optimization Course Project. MIPT}\\[0.35em]
  {\Large Gennady Gotovskii \quad Andrey Gabidulin \quad Leonid Sivachev}\\[0.2em]
  {\large May 2026}
\end{center}

\vspace{0.9em}
{\color{rulegray}\rule{\textwidth}{1.2pt}}
\vspace{0.6em}

\begin{columns}[t,totalwidth=\textwidth]

\begin{column}{0.315\textwidth}
\color{textgray}

\sectiontitle{1. Introduction}
\justifying
Large pretrained vision models are expensive to fully fine-tune: gradients and
optimizer states must be stored for all backbone parameters. Parameter-efficient
fine-tuning methods reduce this cost by training only a small subset of
parameters.

LoRA is one of the most practical methods for adapting Transformer models, but
it can still underperform full fine-tuning. We study whether this gap is partly
explained by gradient mismatch in low-rank adaptation.

\sectiontitle{2. Problem Statement}
We consider supervised image classification:
\[
  \min_{\Phi} \frac{1}{N}\sum_{i=1}^{N}
  \ell(f_{\Phi}(x_i), y_i).
\]

We compare several adaptation strategies for the same pretrained
CLIP-ViT-B/16 backbone: zero-shot inference, linear probing, LoRA, DoRA and
LoRA-Pro.

\sectiontitle{3. LoRA Baseline}
For a pretrained weight matrix $W_0 \in \mathbb{R}^{d \times k}$, LoRA uses
\[
  W = W_0 + sBA,\qquad
  B\in\mathbb{R}^{d\times r},\quad
  A\in\mathbb{R}^{r\times k},\quad
  r\ll \min(d,k).
\]

LoRA freezes $W_0$ and trains only $A$ and $B$. This reduces the number of
trainable parameters from $dk$ to $r(d+k)$ while preserving the original
pretrained backbone.

In CLIP-ViT, adapters are inserted into attention projection matrices, primarily
\code{q\_proj} and \code{v\_proj}. The original projection
\[
  h = W_0x
\]
is replaced by
\[
  h = W_0x + sBAx.
\]

\sectiontitle{Figure 1. LoRA Adapter Architecture}

\begin{center}
\begin{tikzpicture}[
    scale=0.86,
    transform shape,
    >=Latex,
    font=\large,
    input/.style={
        rectangle,
        draw=rulegray,
        line width=1.1pt,
        rounded corners=4pt,
        minimum width=5.0cm,
        minimum height=0.95cm,
        align=center,
        fill=white
    },
    frozen/.style={
        rectangle,
        draw=titleblue,
        line width=1.3pt,
        rounded corners=6pt,
        minimum width=6.4cm,
        minimum height=1.9cm,
        align=center,
        fill=blue!8
    },
    trainable/.style={
        rectangle,
        draw=orange!80!black,
        line width=1.3pt,
        rounded corners=6pt,
        minimum width=5.8cm,
        minimum height=1.45cm,
        align=center,
        fill=orange!15
    },
    output/.style={
        rectangle,
        draw=rulegray,
        line width=1.1pt,
        rounded corners=4pt,
        minimum width=5.2cm,
        minimum height=0.95cm,
        align=center,
        fill=white
    },
    arrow/.style={
        ->,
        line width=1.4pt,
        draw=titleblue
    },
    loraarrow/.style={
        ->,
        line width=1.4pt,
        draw=orange!85!black
    }
]

\node[input] (x) at (0,-2) {$x$};

\node[frozen] (w0) at (-5.6,4.1) {
    \textbf{Frozen}\\[-0.1em]
    pretrained weight\\
    $W_0$
};

\node[trainable] (A) at (5.6,1.0) {
    \textbf{Trainable}\\[-0.1em]
    $A \in \mathbb{R}^{r\times k}$
};

\node[trainable] (B) at (5.6,6.35) {
    \textbf{Trainable}\\[-0.1em]
    $B \in \mathbb{R}^{d\times r}$
};

\node[output] (sum) at (0,10.05) {$+$};

\node[output] (h) at (0,12.85) {
    $h = W_0x + sBAx$
};

\draw[arrow] (x) -- (w0);
\draw[arrow] (w0) -- node[above left, text=mutedgray, pos=0.45] {$W_0x$} (sum);

\draw[loraarrow] (x) -- (A);
\draw[loraarrow] (A) -- node[right, text=mutedgray, pos=0.5] {$Ax$} (B);
\draw[loraarrow] (B) -- node[above right, text=mutedgray, pos=0.45] {$sBAx$} (sum);

\draw[arrow] (sum) -- (h);

\draw[mutedgray, dashed, line width=0.9pt]
    ($(A.east)+(0.35,0)$) -- ($(B.east)+(0.35,0)$)
    node[midway, right=0.25cm, align=left, text=mutedgray, font=\normalsize]
    {$r$-dimensional\\bottleneck};

\node[align=left, text=mutedgray, font=\normalsize] at (0,-6.65) {
    \textcolor{titleblue}{blue branch}: frozen pretrained projection\\
    \textcolor{orange!85!black}{orange branch}: trainable LoRA adapter
};

\end{tikzpicture}
\end{center}

{\normalsize
LoRA keeps the pretrained projection frozen and adds a trainable low-rank branch.
The input is compressed to an $r$-dimensional bottleneck by $A$, expanded by $B$,
scaled by $s=\alpha/r$, and added to the original projection.
}
\end{column}

\begin{column}{0.315\textwidth}
\color{textgray}

\sectiontitle{4. LoRA-Pro: Gradient Mismatch}
Standard LoRA does not directly follow the full fine-tuning gradient. Gradients
over $A$ and $B$ induce an equivalent update in the full weight space:
\[
  dW = s\,dB\,A + s\,B\,dA.
\]

The induced full-space direction is
\[
  \widetilde g = sB g_A + s g_B A,
  \qquad
  g=\partial\mathcal{L}/\partial W.
\]
LoRA-Pro chooses corrected low-rank gradients so that \(\widetilde g\) is closer
to the full fine-tuning gradient:
\[
  \min_{g_A,g_B}\|sBg_A + sg_BA - g\|_F^2.
\]

\sectiontitle{Figure 2. Gradient Mismatch Intuition}

\begin{center}
\begin{tikzpicture}[
    >=Latex,
    scale=1.23,
    transform shape,
    font=\small,
    vector/.style={
        ->,
        line width=1.25pt
    },
    fullgrad/.style={
        vector,
        draw=doragreen
    },
    ga/.style={
        vector,
        draw=titleblue!65!black
    },
    gb/.style={
        vector,
        draw=orange!85!black
    },
    sumgrad/.style={
        vector,
        draw=mutedgray,
        dashed
    },
    mismatch/.style={
        <->,
        line width=0.9pt,
        dashed
    }
]

\draw[
    draw=rulegray,
    fill=panelgray,
    line width=0.8pt,
    opacity=0.75
] (0,0) ellipse (5.2cm and 1.25cm);


\coordinate (o) at (-4.35,-0.35);
\coordinate (gaend) at (-3.75,0.68);
\coordinate (gbend) at (1.25,-0.3);
\coordinate (sum) at ($(gaend)+(gbend)-(o)$);
\coordinate (g) at ($(sum)+(-0.5,4.55)$);

\draw[ga] (o) -- (gaend)
    node[pos=0.54, left, text=titleblue] {$sBg_A$};

\draw[gb] (o) -- (gbend)
    node[pos=0.58, below, text=orange!85!black] {$sg_BA$};

\draw[sumgrad] (o) -- (sum)
    node[pos=1.3, below, text=mutedgray] {$\widetilde g=sBg_A+sg_BA$};

\fill[mutedgray] (sum) circle (1.5pt);

\draw[fullgrad] (o) -- (g)
    node[pos=0.53, left, text=doragreen] {$g=\partial\mathcal{L}/\partial W$};

\draw[mismatch, draw=mutedgray] (sum) -- (g)
    node[pos=0.45, left, text=mutedgray] {\scriptsize mismatch};

\node[
    rectangle,
    rounded corners=4pt,
    draw=orange!85!black,
    fill=orange!8,
    line width=0.9pt,
    align=center,
    minimum width=5.5cm,
    minimum height=1.05cm
] at (7.15,3.05)
{
    {\scriptsize LoRA-Pro gradient correction}\\[-0.15em]
    \(\displaystyle
    \min_{g_A,g_B}
    \|sBg_A + sg_BA - g\|_F^2
    \)
};

\node[align=left, text=mutedgray] at (1.5,-4.05) {
    \textcolor{doragreen}{green}: full fine-tuning gradient \(g\)\\
    \textcolor{titleblue!65!black}{dark blue}: contribution \(sBg_A\)\\
    \textcolor{orange!85!black}{orange}: contribution \(sg_BA\)\\
    \textcolor{mutedgray}{gray dashed}: equivalent gradient \(\widetilde g\)
};

\end{tikzpicture}
\end{center}

{\small
The gradients \(g_A\) and \(g_B\) induce an equivalent full-weight gradient
\(\widetilde g=sBg_A+sg_BA\) inside the reachable family. LoRA-Pro adjusts
\(g_A\) and \(g_B\) so that this induced gradient moves closer to the full
fine-tuning gradient \(g\).
}

\sectiontitle{5. Method Comparison}
\begin{center}
\begin{tabular}{>{\bfseries}p{0.22\linewidth}p{0.66\linewidth}}
\toprule
Method & Adaptation strategy \\
\midrule
Zero-shot & No training; compare image and text embeddings. \\
Linear probe & Freeze CLIP visual encoder; train classifier head. \\
LoRA & Train low-rank adapters in attention projections. \\
DoRA & Decompose weights into magnitude and direction; apply LoRA-style updates to direction. \\
LoRA-Pro & Correct low-rank gradients to reduce mismatch. \\
\bottomrule
\end{tabular}
\end{center}

\sectiontitle{6. Experimental Setup}
\begin{center}
\begin{tabular}{>{\bfseries}ll}
\toprule
Backbone & CLIP-ViT-B/16 \\
Checkpoint & \code{openai/clip-vit-base-patch16} \\
Dataset & EuroSAT first, CIFAR-100 second \\
Rank & $r=8$ \\
Target modules & \code{q\_proj}, \code{v\_proj} \\
Metrics & accuracy, trainable params, time, memory, gradient discrepancy \\
\bottomrule
\end{tabular}
\end{center}

\end{column}

\begin{column}{0.315\textwidth}
\color{textgray}

\sectiontitle{7. Preliminary Experiments}
We ran initial EuroSAT baselines using CLIP-ViT-B/16. Zero-shot CLIP gives a
sanity-check baseline, linear probing measures frozen-feature quality, and
standard LoRA with rank $r=8$ gives the first trainable PEFT baseline.

\begin{center}
\begin{tabular}{p{0.30\linewidth}p{0.16\linewidth}p{0.18\linewidth}p{0.25\linewidth}}
\toprule
Method & Dataset & Accuracy & Params \\
\midrule
Zero-shot CLIP & EuroSAT & \textbf{49.59\%} & 0 \\
Linear Probe & EuroSAT & \textbf{93.26\%} & 5.1K \\
LoRA $r=8$ & EuroSAT & \textbf{97.81\%} & 496.7K \\
\midrule
DoRA $r=8$ & EuroSAT & next & adapter + magnitude \\
LoRA-Pro $r=8$ & EuroSAT & next & adapter + correction \\
\bottomrule
\end{tabular}
\end{center}

{\small
LoRA trains \(496{,}650\) of \(150{,}117{,}387\) parameters, approximately
\(0.33\%\) of the model, while reaching \(97.81\%\) validation accuracy. These
results validate the pipeline and establish the baseline for LoRA-Pro.
}

\sectiontitle{8. Expected Results and Next Steps}
The preliminary results confirm that the GPU training and evaluation pipeline
works. Next, we compare LoRA, DoRA and LoRA-Pro under the same rank and target
modules.

The main question is whether LoRA-Pro improves optimization dynamics by reducing
normalized gradient mismatch:
\[
  \Delta_{\mathrm{grad}}=\frac{\|\widetilde g-g\|_F}{\|g\|_F}.
\]
We will report accuracy, trainable parameters, time, memory and gradient
mismatch.

\sectiontitle{Figure 3. EuroSAT Baseline Accuracy}

\begin{center}
\begin{tikzpicture}[
    x=0.86cm,
    y=0.075cm,
    font=\normalsize,
    bar/.style={
        rounded corners=3pt
    }
]

% Axes
\draw[->, line width=1.1pt, draw=mutedgray] (0,0) -- (11.2,0);
\draw[->, line width=1.1pt, draw=mutedgray] (0,0) -- (0,104);

% Y ticks
\foreach \y in {0,20,40,60,80,100} {
    \draw[draw=rulegray, line width=0.7pt] (-0.12,\y) -- (0,\y);
    \node[left, text=mutedgray, font=\normalsize] at (-0.18,\y) {\y};
    \draw[draw=rulegray, line width=0.35pt, opacity=0.45] (0,\y) -- (10.6,\y);
}

% Bars
\draw[bar, fill=zeroshotgray!70, draw=zeroshotgray] (1.0,0) rectangle (2.6,49.59);
\draw[bar, fill=linearblue!75, draw=linearblue] (4.4,0) rectangle (6.0,93.26);
\draw[bar, fill=lorapurple!75, draw=lorapurple] (7.8,0) rectangle (9.4,97.81);

% Labels
\node[below, align=center, text=textgray, font=\small] at (1.8,-7) {Zero-shot\\CLIP};
\node[below, align=center, text=textgray, font=\small] at (5.2,-7) {Linear\\Probe};
\node[below, align=center, text=textgray, font=\small] at (8.6,-7) {LoRA\\$r=8$};

\node[above, text=zeroshotgray!80!black, font=\bfseries\small] at (1.8,51.5) {49.59\%};
\node[above, text=linearblue!80!black, font=\bfseries\small] at (5.2,95.2) {93.26\%};
\node[above, text=lorapurple!80!black, font=\bfseries\small] at (8.6,99.5) {97.81\%};

\node[rotate=90, text=mutedgray, font=\small] at (-1.65,52) {Accuracy (\%)};

\end{tikzpicture}
\end{center}

{\small
CLIP features are already strong on EuroSAT, while LoRA further improves
accuracy with less than \(0.4\%\) trainable parameters.
}

\sectiontitle{9. Publication Plan and Conclusion}
\begin{itemize}
  \item Mid-May 2026: poster draft and completed EuroSAT baselines.
  \item Late May 2026: complete EuroSAT and CIFAR-100 runs for LoRA, DoRA, and LoRA-Pro.
  \item Early June 2026: finalize the report with reproducible code, configs, and result analysis.
  \item June 2026: publish the final materials on GitHub or a small project page.
\end{itemize}

This project tests whether gradient mismatch is a practical bottleneck for
low-rank adaptation of vision foundation models. The completed baselines show
that the pipeline is reproducible and that LoRA is a strong starting point for
the LoRA-Pro comparison.

\end{column}

\end{columns}

\end{frame}
\end{document}
